% Generated draft: requires mathematical and source review.
\documentclass[11pt]{article}
\usepackage[margin=25mm]{geometry}
\usepackage{fontspec}
\setmainfont{Latin Modern Roman}
\usepackage{amsmath,amssymb,mathtools,unicode-math}
\setmathfont{Latin Modern Math}
\usepackage{enumitem,needspace,xurl,hyperref}
\setlength{\parindent}{0pt}
\setlength{\parskip}{5pt}
\setlength{\emergencystretch}{3em}
\newcommand{\N}{\mathbb{N}}
\newcommand{\Z}{\mathbb{Z}}
\newcommand{\Q}{\mathbb{Q}}
\newcommand{\R}{\mathbb{R}}
\newcommand{\C}{\mathbb{C}}
\newcommand{\F}{\mathbb{F}}
\newcommand{\A}{\mathbb{A}}
\newcommand{\PP}{\mathbb{P}}
\newcommand{\E}{\mathbb{E}}
\newcommand{\eps}{\varepsilon}
\newcommand{\dd}{\,\mathrm{d}}
\begin{document}
\section*{UnsolvedMath ID: 10400209}
\textbf{Source problem code:} catalogue entry
\subsection{Definitions and mathematical statement}
Let \(D\) be a diagram of a knot \(K\subset S^3\). In the intended knot-group interpretation, put
\[
 G_K=\pi_1(S^3\setminus K)
 =\langle x_1,\ldots,x_m\mid r_1,\ldots,r_s\rangle,
\]
using a Wirtinger presentation from \(D\), with the \(x_i\) representing consistently oriented meridians. A homomorphism to a finite permutation group is made explicit by giving permutations \(\rho(x_i)\) satisfying every relation \(r_j\).

\emph{Problem (H.R. Morton).} From \(D\), explicitly find a homomorphism of the kind intended by the source to some permutation group, or establish that \(K\) is the trivial knot. The phrase ``such homomorphism'' refers to preceding material absent from the supplied record. In particular, a condition distinguishing it from an arbitrary trivial or cyclic-image homomorphism is missing. A noncyclic-image homomorphism would be a natural nontriviality certificate in the knot-group interpretation, but the excerpt does not establish that this is exactly the intended condition; no particular admissibility criterion is silently substituted.

The requested refinements are:
\begin{enumerate}
\item Give an upper bound, in terms of \(D\), for the orders of the permutation groups that need be considered in this search. The source asks about group order, not merely permutation degree; a precise formulation must also distinguish the target group from the image subgroup when relevant.
\item Determine what happens when the images of meridians, which are all conjugate, are required to have a specified cycle type; the source gives a single transposition as an example. In an embedding of the target into \(S_n\), this is a restriction to a chosen conjugacy class of permutations, for example type \((2,1^{n-2})\). The intended allowed degrees and the criterion for a successful restricted search are not supplied.
\end{enumerate}
The input-diagram construction and both refinements are retained without claiming that an arbitrary homomorphism answers the problem or that every chosen cycle type must detect every nontrivial knot.
\subsection{Short English statement}
Morton asks for an explicit finite permutation representation from a knot diagram, or a proof that the knot is trivial, with a diagram-based bound on target-group orders and an investigation of prescribed meridian cycle types. The essential condition referred to by ``such homomorphism'' is missing from the excerpt.
\subsection{Sources}
\begin{itemize}
\item Local UnsolvedMath registry, permanent ID 10400209, source code catalogue entry; attribution H.R. Morton. The stored title is truncated after ``Problem 12.3 --- (H.R.''.
\item Ohtsuki, \emph{Problems on invariants of knots and 3-manifolds} (2002), Problem 12.3, PDF page 159. Supplied URL: \url{https://msp.org/gtm/2002/04/gtm-2002-04-024s.pdf}; recorded access 2026-07-29; extraction pdf-font-structure. The collection distinguishes 231 numbered bold problem, conjecture, and question headings.
\item The supplied registry labels the item open and OPEN-TRIAGE and describes explicit finite-group homomorphisms from knot diagrams. Current primary-literature status and redistribution terms for the publicly accessible author/publisher copy remain NEEDS\_REVIEW.
\end{itemize}
\end{document}
